\documentclass[11pt]{article}
\usepackage[a4paper,margin=2.5cm]{geometry}
\usepackage[T1]{fontenc}
\usepackage[utf8]{inputenc}
\usepackage{lmodern}
\usepackage{amsmath,amssymb,booktabs,array}
\usepackage{siunitx}
\usepackage{enumitem}
\usepackage{hyperref}
\usepackage{xcolor}

\sisetup{
  round-mode=places,
  round-precision=3,
  scientific-notation=true,
  table-number-alignment=center
}

\title{Technical Analysis of the Experimental Results\\for the Practical Grover--Rudolph Implementation}
\author{Prepared from the files\\\texttt{gr\_log\_20260308-131538.csv}, \texttt{gr\_summary\_20260308-131538.csv},\\and the readout calibration matrix \texttt{Mfull\_20260308-112738.npy}}
\date{March 8, 2026}

\begin{document}
\maketitle

\section*{Executive summary}

The results provide a clear and coherent picture.

\begin{enumerate}[label=\arabic*.]
\item The \textbf{ideal simulation is essentially exact}. The row \emph{Target vs SIM ideal (FULL)} reports TV and $\ell^2$ errors at numerical precision and fidelity equal to $1$. Therefore, the algorithmic construction and ideal compilation are correct.

\item The \textbf{hardware execution is viable but still noticeably degraded}. For the full circuit, the average NMR result reaches a classical fidelity of about $0.927$ with respect to the target distribution, but the total variation distance remains around $0.245$.

\item \textbf{Readout mitigation does not improve the final full-circuit result}. In fact, the mitigated average is slightly worse than the raw average at the FULL stage. This strongly suggests that the dominant error source is no longer measurement bias alone, but accumulated gate errors, control imperfections, and/or decoherence in the deeper part of the circuit.

\item The staged execution \textbf{localizes the main degradation to the last layer}. The first stage is reasonably close to its stage-ideal output, the intermediate stage is still acceptable, and the largest deterioration appears when the final controlled-rotation block is added.
\end{enumerate}

\section{Data used in the analysis}

The analysis is based on the uploaded output files:
\begin{itemize}
\item \texttt{gr\_summary\_20260308-131538.csv}
\item \texttt{gr\_log\_20260308-131538.csv}
\item \texttt{Mfull\_20260308-112738.npy}
\end{itemize}

The summary file reports the global FULL-stage comparison:
\begin{center}
\begin{tabular}{lccc}
\toprule
Comparison & TV & $\ell^2$ & Fidelity \\
\midrule
Target vs SIM ideal (FULL) & 1.380e-10 & 1.024e-10 & 1.000000 \\
Target vs NMR raw avg (FULL) & 0.245245 & 0.192330 & 0.927042 \\
Target vs NMR mitigated (FULL, ridge=0.001) & 0.250522 & 0.202646 & 0.922016 \\
Target vs NMR mixed (FULL, lam=0.3) & 0.246828 & 0.195013 & 0.925782 \\
\bottomrule
\end{tabular}
\end{center}

The readout calibration matrix is numerically well behaved:
\[
\operatorname{cond}(M_{\mathrm{full}})\approx 1.369,
\qquad
\text{mean diagonal entry}\approx 0.855.
\]
Hence the inversion problem is not ill-conditioned, and the failure of mitigation at the FULL stage should not be attributed to a pathological calibration matrix.

\section{Main findings}

\subsection{Ideal correctness is fully validated}

The ideal simulation reproduces the target distribution with errors of order $10^{-10}$. This is decisive: the discrepancy observed experimentally is not caused by a theoretical or software-level mistake in the Grover--Rudolph construction itself. The implementation is correct in the noiseless model.

\subsection{The full hardware circuit is meaningful but not yet precise}

For the FULL circuit, the relevant averages are:
\[
\mathrm{TV}_{\mathrm{raw}} \approx 0.245,
\qquad
\mathrm{TV}_{\mathrm{mitig}} \approx 0.251,
\]
\[
F_{\mathrm{raw}} \approx 0.927,
\qquad
F_{\mathrm{mitig}} \approx 0.922.
\]

This means that the experimental preparation retains a substantial part of the intended probability profile, but still redistributes roughly one quarter of the probability mass with respect to the target. In practical terms, the experiment is \emph{successful as a proof of viability} but \emph{not yet accurate as a high-fidelity state-preparation primitive}.

\subsection{Readout mitigation is not the dominant lever at the final stage}

A particularly important observation is that mitigation helps little---and actually worsens the final FULL average slightly:
\[
\mathrm{TV}:~ 0.245245 \to 0.250522,
\qquad
F:~ 0.927042 \to 0.922016.
\]

This indicates that once the full ladder has been executed, the dominant error budget is likely coherent or dynamical rather than purely classical measurement bias. In other words, correcting the readout layer is no longer enough because the state arriving at measurement is already significantly distorted.

\section{Stage-by-stage diagnosis}

For the staged experiments, the most informative quantities are the distances to the stage-specific simulation outputs, namely \texttt{tv\_vs\_sim} and \texttt{fid\_vs\_sim}. The averages over the six NMR repetitions are:

\begin{center}
\begin{tabular}{lccccc}
\toprule
Stage & Runs & Mean TV vs sim & Std TV vs sim & Mean fidelity vs sim & Std fidelity vs sim \\
\midrule
L0 & 6 & 0.085 & 0.019 & 0.915 & 0.020 \\
L01 & 6 & 0.119 & 0.017 & 0.884 & 0.018 \\
FULL & 6 & 0.256 & 0.036 & 0.915 & 0.027 \\
\bottomrule
\end{tabular}
\end{center}

The pattern is structurally clear:
\begin{itemize}
\item \textbf{L0} is already reasonably close to its intended stage output.
\item \textbf{L01} shows additional degradation, but still remains in a comparable regime.
\item \textbf{FULL} exhibits the largest cumulative deviation once the last control structure is appended.
\end{itemize}

This supports the interpretation that the final uniformly controlled rotation block, or its decomposition into native operations, is the principal experimental bottleneck.

\section{Distribution-level interpretation at the FULL stage}

The target full probability vector is
\[
p_{\mathrm{target}} =
[0.05,\,0.10,\,0.15,\,0.20,\,0.20,\,0.15,\,0.10,\,0.05].
\]

The average experimental vectors are:
\[
p_{\mathrm{raw}} \approx [0.1068, 0.1783, 0.1412, 0.1171, 0.0878, 0.1087, 0.1675, 0.0926],
\]
\[
p_{\mathrm{mitig}} \approx [0.0957, 0.1765, 0.1454, 0.1103, 0.0829, 0.1109, 0.1923, 0.0861].
\]

Two features stand out.

First, the output is not simply flattened by uniform noise. Instead, some basis states are clearly overweighted and others underweighted, which is more compatible with angle miscalibration, imperfect conditional rotations, residual coupling effects, or accumulated two-qubit gate errors than with a purely depolarizing picture.

Second, the mitigated average does not move the distribution closer to the target in a meaningful way. It redistributes weight, but not in the directions needed to recover the intended final profile.

\section{Marginal biases}

For the mitigated FULL average, the one-qubit marginals are approximately
\[
P(q_0=1)\approx 0.472,\qquad
P(q_1=1)\approx 0.534,\qquad
P(q_2=1)\approx 0.484.
\]

In the target full distribution, all three marginals should be exactly balanced:
\[
P(q_j=1)=0.5,\qquad j=0,1,2.
\]

The observed deviations are moderate but noticeable, especially on $q_0$ and $q_1$. This again points to a control imbalance rather than a purely symmetric readout distortion.

\section{Run-to-run variability}

At the FULL stage, the six NMR repetitions show a non-negligible spread:
\[
\mathrm{TV}_{\mathrm{target}} \in [0.212, 0.310],
\qquad
F_{\mathrm{target}} \in [0.871, 0.948].
\]

This is not catastrophic, but it is large enough to suggest temporal drift or sensitivity to calibration conditions across runs. For an NMR platform, such variability is entirely plausible and should be explicitly tracked in any subsequent campaign.

\section{Overall conclusion}

The experimental campaign supports the following technical conclusion:

\begin{quote}
The practical Grover--Rudolph implementation is \textbf{algorithmically correct} and \textbf{experimentally viable} on the 3-qubit NMR platform, but the dominant limitation at the FULL stage is \textbf{not readout}. The principal bottleneck is the \textbf{accumulation of circuit-level errors}, most likely concentrated in the last controlled-rotation layer and its native-gate decomposition.
\end{quote}

A fair and precise summary would be:
\begin{itemize}
\item \textbf{Validated}: ideal implementation and staged execution logic.
\item \textbf{Encouraging}: full-circuit hardware fidelity remains above $0.92$ on average.
\item \textbf{Limiting factor}: TV error around $0.25$ is still too large for precise state preparation.
\item \textbf{Key diagnostic}: readout mitigation does not rescue the final output, so future improvements should focus primarily on pulse-level and gate-level optimization.
\end{itemize}

\section{Recommended next steps}

The most useful next experimental actions are the following:
\begin{enumerate}[label=\arabic*.]
\item \textbf{Isolate the final layer.} Benchmark the last controlled-rotation block separately, including its native decomposition depth, pulse duration, and sensitivity to calibration drift.
\item \textbf{Compare raw and mitigated stage outputs systematically.} The current evidence suggests that mitigation is useful only for shallow stages. That claim should be documented explicitly in a dedicated table.
\item \textbf{Track temporal drift.} Group runs by acquisition time and re-estimate whether the distribution changes during the campaign.
\item \textbf{Reduce circuit-level sensitivity.} Explore alternative decompositions of the final UCRy block, gate cancellations, shorter pulse schedules, or angle clipping when physically justified.
\item \textbf{Report uncertainty bars.} For future figures and papers, include mean $\pm$ standard deviation over repeated runs for TV, $\ell^2$, and fidelity.
\end{enumerate}

\section*{Suggested wording for a paper or report}

A concise statement suitable for inclusion in a manuscript is:

\begin{quote}
The staged NMR experiments confirm that the Grover--Rudolph state-preparation protocol is implemented correctly in simulation and remains viable on hardware. However, the final FULL circuit exhibits a total variation error of about $0.25$ despite a classical fidelity above $0.92$. Since readout mitigation does not improve the final result, the dominant experimental limitation is attributed to accumulated circuit-level errors, most likely in the last controlled-rotation layer rather than in the measurement stage alone.
\end{quote}

\end{document}
